% sec:kerr-wormhole — The Morris-Thorne Wormhole
\section{The Morris--Thorne Wormhole\starmark{}}
\label{sec:kerr-wormhole}
\coderef{Spacetime}{MorrisThorne}

The Morris--Thorne wormhole is a static, spherically symmetric
solution to the Einstein field equations that replaces the
singularity and event horizon with a smooth throat connecting two
asymptotically flat regions \cite{MorrisThorne1988}.  No photon is
ever captured: every geodesic either passes through the throat to the
other universe or deflects back to infinity.  The bright lensing ring
formed by unstable photon orbits at the throat, together with the
double sky visible through it, produces images visually distinct from
any black hole --- making the wormhole an ideal stress test for the
geodesic integrator.

\paragraph{The line element.}

The zero-tidal-force Morris--Thorne metric, written in terms of the
proper radial distance $\ell$ from the throat, is
%
\begin{equation}
  \d s^2 = -\d t^2 + \d\ell^2
    + \bigl(b_0^2 + \ell^2\bigr)
      \bigl(\d\theta^2 + \sin^2\!\theta\;\d\varphi^2\bigr) \,,
  \label{eq:mt-metric}
\end{equation}
%
reproducing \cref{eq:morris-thorne-metric} for convenience.  The
parameter $b_0$ is the throat radius: the circumferential radius
$r(\ell) = \sqrt{b_0^2 + \ell^2}$ reaches its minimum $b_0$ at
$\ell = 0$ and grows without bound as $\abs{\ell} \to \infty$.
The coordinate $\ell$ runs from $-\infty$ to $+\infty$, with
negative values labelling one asymptotic region and positive values
the other.

Three properties set this geometry apart from the Schwarzschild
and Kerr spacetimes.  First, $g_{tt} = -1$ is everywhere negative
--- the time Killing vector never changes character, so there is no
ergoregion and no horizon.
\physnote{Setting $g_{tt} = -1$ everywhere means there is no
gravitational redshift: a photon emitted at the throat arrives at
infinity with unchanged frequency.  The general Morris--Thorne
family includes a redshift function $\Phi(\ell)$ in
$g_{tt} = -e^{2\Phi}$, but the codebase implements only the
$\Phi = 0$ case.}
Second, the angular part
$(b_0^2 + \ell^2)(\d\theta^2 + \sin^2\!\theta\;\d\varphi^2)$ is a
sphere of radius $r(\ell)$ that shrinks to a minimum at $\ell = 0$
but never pinches off --- the topology is
$\mathbb{R} \times S^2$, not $\mathbb{R}^3$.  Third, the metric
coefficients are everywhere finite and smooth, so there is no
coordinate or curvature singularity anywhere.

\paragraph{Cartesian form and the codebase.}

The geodesic integrator works in Cartesian coordinates, not the
$(\ell, \theta, \varphi)$ chart of \cref{eq:mt-metric}.  The
conversion identifies the Euclidean distance
$\rho = \sqrt{x^2 + y^2 + z^2}$ with the magnitude of the proper
radial distance, $\rho = \abs{\ell}$, so that the throat maps to
the Cartesian origin $\rho = 0$.  Writing $n_i = x_i/\rho$ for the
unit radial vector, the flat-space line element decomposes as
$\d\rho^2 = n_i n_j \,\d x^i \d x^j$ for the radial part and
$\rho^2(\d\theta^2 + \sin^2\!\theta\;\d\varphi^2) =
  (\delta_{ij} - n_i n_j)\,\d x^i \d x^j$
for the angular part.  Since $\d\ell^2 = \d\rho^2$ and
$b_0^2 + \ell^2 = b_0^2 + \rho^2$, substituting into
\cref{eq:mt-metric} gives the spatial metric
%
\begin{equation}
  g_{ij} = \Bigl(1 + \frac{b_0^2}{\rho^2}\Bigr)\delta_{ij}
           - \frac{b_0^2}{\rho^2}\,n_i\,n_j \,,
  \label{eq:mt-cartesian-spatial}
\end{equation}
%
with $g_{tt} = -1$ and $g_{ti} = 0$.  The eigenvalues of $g_{ij}$
separate cleanly: the radial eigenvalue is
$(1 + b_0^2/\rho^2) - b_0^2/\rho^2 = 1$, confirming that $\rho$
measures proper radial distance, and each transverse eigenvalue is
$1 + b_0^2/\rho^2 = r^2/\rho^2$, encoding the ratio of the
circumferential radius $r = \sqrt{b_0^2 + \rho^2}$ to the coordinate
radius.  Near the throat this ratio diverges: a circle at small
$\rho$ has circumference $2\pi r \gg 2\pi\rho$.
\warnnote{The Cartesian chart covers a single sheet of the wormhole
($\ell \geq 0$, say) and maps the throat to the origin $\rho = 0$.
The metric coefficients diverge as $\rho \to 0$ --- a coordinate
singularity, not a physical one.  The codebase handles throat
crossings in the Rust integrator, which clamps $\rho$ near the
throat and matches the geodesic onto the second sheet.}
\implnote{The function \texttt{morrisThorneMetricFn} implements
\cref{eq:mt-cartesian-spatial} directly: it takes the throat radius
$b_0$ and a coordinate list $[t, x, y, z]$, returning the ten
independent metric components in \texttt{Sym4} order.  The
polymorphic type signature
\texttt{Floating a => a -> [a] -> [a]} enables automatic
differentiation of the Christoffel symbols, following the same
strategy as the Kerr metric (\cref{sec:kerr-ks}).}

The inverse spatial metric follows from the eigenvalue decomposition:
inverting $1$ in the radial direction and $\rho^2/(\rho^2 + b_0^2)$
in each transverse direction gives
%
\begin{equation}
  g^{ij} = \frac{\rho^2}{\rho^2 + b_0^2}\,\delta^{ij}
           + \frac{b_0^2}{\rho^2 + b_0^2}\,n^i\,n^j \,.
  \label{eq:mt-cartesian-inverse}
\end{equation}
%
Unlike the Kerr--Schild inverse, which follows from a rank-one null
perturbation (\cref{sec:kerr-ks}), the Morris--Thorne inverse relies
on the radial--transverse decomposition of $g_{ij}$, making the
inversion algebraic rather than requiring a general $3\times 3$
matrix calculation.

\paragraph{Throat geometry.}

The throat at $\ell = 0$ is a two-sphere of circumference
$2\pi b_0$.  In the $(\ell, \theta, \varphi)$ chart, the metric
\cref{eq:mt-metric} is perfectly smooth there: $g_{\ell\ell} = 1$
and $g_{\theta\theta} = b_0^2$, so the throat is a regular minimal
surface rather than a coordinate singularity.  Every radial geodesic
passes through this minimal sphere and emerges in the other
asymptotic region.  The embedding diagram of
\cref{fig:embedding-diagrams} shows the geometry as a smooth funnel,
symmetric about $\ell = 0$, whose circumferential radius
$r(\ell) = \sqrt{b_0^2 + \ell^2}$ increases monotonically away from
the throat.
\physnote{The curvature is concentrated near the throat.  Far
from it ($\abs{\ell} \gg b_0$, equivalently $\rho \gg b_0$),
$b_0^2/\rho^2 \to 0$ and the metric approaches flat Minkowski space,
as required for an asymptotically flat spacetime.  The Riemann tensor
components decay as $b_0^2/r^4$.}

\paragraph{Geodesic structure.}

The effective potential for null geodesics in the equatorial plane
follows from the constants of motion.  With conserved energy $E$ and
angular momentum $L$, the radial equation of motion reduces to
%
\begin{equation}
  \dot{\ell}^{\,2} = E^2 - \frac{L^2}{b_0^2 + \ell^2} \,,
  \label{eq:mt-null-radial}
\end{equation}
%
where the dot denotes differentiation with respect to an affine
parameter.  The effective potential
$V(\ell) = L^2/(b_0^2 + \ell^2)$ has a unique maximum at
$\ell = 0$: the throat is an unstable circular photon orbit, the
wormhole's analogue of the black hole photon sphere
(\cref{sec:schw-orbits}).

The critical impact parameter follows immediately: setting
$\dot{\ell} = 0$ at $\ell = 0$ gives $b_c \equiv L/E = b_0$.
Photons with $b < b_0$ have enough energy to cross the potential
barrier and thread the throat to the other universe.  Photons with
$b > b_0$ are reflected by the barrier and return to infinity.
Photons with $b \approx b_0$ wind multiple times around the throat
before escaping, producing a bright lensing ring in the rendered
image --- formed by the same mechanism of unstable circular orbits
that creates the photon ring of a black hole
(\cref{sec:schw-orbits}).

Inside this ring, the observer sees a second copy of the sky: the
universe on the other side of the throat.  The absence of an event
horizon means no photons are absorbed, so the wormhole image has no
shadow --- a stark visual contrast with the dark Kerr silhouette of
\cref{sec:kerr-horizons}.

\paragraph{Exotic matter and the NEC.}

The price of a traversable throat is exotic matter.  The Einstein
field equations require the stress-energy tensor sustaining the
throat to violate the null energy condition
$\stressenergy\, k^\mu k^\nu \geq 0$ for all null $k^\mu$
(\cref{sec:gr-stress-energy}).  No known classical matter
achieves this violation in a controlled setting
\cite{MorrisThorne1988}.  The Morris--Thorne spacetime is therefore a
mathematical existence proof rather than an astrophysical prediction
--- but as a test bed for the ray tracer, its value lies precisely in
its exotic character: the absence of horizons, the nontrivial
topology, and the double-sky images all exercise code paths that
black hole spacetimes leave untouched.
